Zunächst konvertieren wir 1 kg in GeV. In natürlichen Einheiten gilt $\hbar = c = 1$, was bedeutet, dass Energie, Masse und Impuls dieselben Einheiten haben. Die Umrechnung von kg in GeV erfolgt über die Ruheenergie eines Kilogramms, $E = mc^2$. Da $c = 1$ in natürlichen Einheiten, ist $c^2$ ebenfalls 1, und wir verwenden die Umrechnung von Joule in GeV:

\[
1 \, \text{kg} = \frac{1 \, \text{kg} \cdot (3 \times 10^8 \, \text{m/s})^2}{1.60218 \times 10^{-10} \, \text{J/GeV}}
\]

Da $1 \, \text{J} = 1 \, \text{kg} \cdot \text{m}^2/\text{s}^2$, erhalten wir:

\[
1 \, \text{kg} \approx 5.61 \times 10^{26} \, \text{GeV}
\]

Nun konvertieren wir 1 s in $\mathrm{GeV}^{-1}$. In natürlichen Einheiten, wo $c = 1$, ist die Lichtgeschwindigkeit gleich 1, was bedeutet, dass Zeit und Länge dieselben Einheiten haben. Die Umrechnung erfolgt über die Beziehung zwischen Energie und Zeit, die durch die Planck-Konstante $\hbar$ vermittelt wird. Da $\hbar = 6.582 \times 10^{-25} \, \text{GeV s}$, ergibt sich:

\[
1 \, \text{s} = \frac{1}{6.582 \times 10^{-25} \, \text{GeV}}
\]

Jetzt berechnen wir die Gravitationskonstante $G_N$ in natürlichen Einheiten. In SI-Einheiten hat $G_N$ die Dimension $\mathrm{m}^3 \, \mathrm{kg}^{-1} \, \mathrm{s}^{-2}$. Um dies in natürlichen Einheiten umzurechnen, verwenden wir:

\[
G_N = 6.67 \times 10^{-11} \, \mathrm{m}^3 \, \mathrm{kg}^{-1} \, \mathrm{s}^{-2} \times \left( \frac{6.582 \times 10^{-25} \, \text{GeV}}{1} \right)^2 \times \left( \frac{1}{5.61 \times 10^{26} \, \text{GeV}} \right)
\]

Dies ergibt:

\[
G_N \approx 6.71 \times 10^{-39} \, \text{GeV}^{-2}
\]

Schließlich berechnen wir die Planck-Masse $M_{Pl}$:

\[
M_{Pl} = G_N^{-1/2} \approx (6.71 \times 10^{-39} \, \text{GeV}^{-2})^{-1/2} \approx 1.22 \times 10^{19} \, \text{GeV}
\]